\section*{Question 2 - Do we really need complex numbers}
Prove that we can assume w.l.o.g.\ that all amplitudes in a quantum computation are real
numbers!\\
\textbf{Guidance:} Show that any quantum circuit on $n$ qubits that uses $t$ two-qubit gates with
complex coefficients can be simulated exactly by a quantum circuit on $n + 1$ qubits, that uses
at most $t$ three-qubit gates with real coefficients.
Show that the probability of acceptance
remains the same in your construction.
Hint: Remember that we assume that the input for a quantum circuit is in the standard
basis (so, in the base of our induction, we know that the state has only real amplitudes).
Use the extra qubit to decide whether the original state is in the Imaginary or Real part,
and adapt each gate to a gate with one more qubit, such that the new gate works in an
equivalent way to the old one.
More formally, suppose that at some step of the original
circuit the state at that step is $\sum_x(a_x + ib_x)\ket{x}$, the state at that step in the new circuit
should be $\sum_x a_x\ket{0x} + b_x\ket{1x}$.

\subsection*{Solution}
Let $M$ be a map that acts on states and operators as follows:
\begin{align*}
    &M(\ket{c_k}&=\sum_x(a_x+ib_x)\ket{x})=\ket{r_k}&=\sum_x a_x\ket{0x}+b_x\ket{1x}\\
    &M(A+iB)=\begin{pmatrix*}
                 A&-B\\
                 B&A
    \end{pmatrix*}
\end{align*}
To prove that we can assume w.l.o.g that all amplitudes in a quantum computation are real numbers we need to show that for any given circuit $c$ with complex amplitudes there exists an equivalent circuit $r$ with real amplitudes - meaning that $r$ and $c$ have the same probabilities \\

\textbf{Induction on the number of gates:}\\
\textbf{Base case ($k=0$):} The input state is a computational basis state, hence has real amplitudes.
Therefore, by the quantum circuit definition, the circuit's input is in the standard basis and is classical - therefore the input's amplitudes are real hence the inputs of $c$ and $r$ are trivially equivalent (i.e. $c_0=r_0$).\\
\textbf{Induction hypothesis:} Assume that for some $\ket{k}$ in circuit $c$, $M(\ket{k})$ in circuit $r$ gives the same probabilities\\
\textbf{Induction step:} Let the next gate be $U_{k+1}$, then:\\
\begin{align*}
    \ket{c_{k+1}}&=U_{k+1}\ket{c_k}\\
    &=(A+iB)(\sum_x a_x+ib_x)\ket{x}\\
    &=(\sum_x(A_x a_x-B_x b_x)+i(A_x b_x+B_x a_x))\ket{x}\\
    \ket{r_{k+1}}&=G_{k+1}^r\ket{r_k}\\
    &=G_{k+1}^r\sum_x \left(a_x\ket{0x}+b_x\ket{1x}\right)\\
    &=\left(\begin{matrix*}
                A & -B\\
                B & A\\
    \end{matrix*}\right)\left(\sum_x a_x\ket{0x}+b_x\ket{1x}\right)\\
    &=\sum_x \left((A_x a_x-B_x b_x)\ket{0x}+(A_x b_x +B_x a_x)\ket{1x}\right)\\
\end{align*}
From above, we get $M(\ket{c_{k+1}})=\ket{r_{k+1}}$.
The probability to get state $\ket{j}$ in circuit $c$ is:
\begin{equation*}
    \braket{k}{c_{k+1}}=\abs{(A_x a_x-B_x b_x)+i(A_x b_x+B_X a_x)}^2=(A_x a_x-B_x b_x)^2+(A_x b_x+B_X a_x)^2
\end{equation*}
The probability to get state $\ket{j}$ in circuit $r$ is:
\begin{equation*}
    \braket{k}{r_{k+1}}=(A_x a_x-B_x b_x)^2+(A_x b_x+B_X a_x)^2
\end{equation*}
Which results in the same probability.
QED\@.
